Once an interaction object has been defined, it is tempting to design an architecture that emits
it directly.  The difficulty is that semantic requirements become information constraints.  If an
edge is claimed to be invariant under substitutions at its endpoints, then the mechanism choosing
that edge cannot depend on those endpoint values.  If a zero edge is claimed to save computation,
the implementation must actually skip its work.  If an update is claimed to descend an energy,
its memories, state variables, and parameter tying must satisfy the assumptions of the descent
proof.

These constraints are easy to blur because neural architectures combine several roles in one
tensor.  An attention coefficient can select a route, scale a message, and contribute to an output
response.  A pair representation can store persistent state or merely cache a transient feature.
A soft gate can express uncertainty without producing exact support.  The operator framework
separates these roles before asking which implementation is efficient.

The chapter develops two kinds of result.  The first kind is semantic: endpoint-removable context,
strict targetwise support, reciprocity, and state--energy compatibility state what an architecture
must satisfy for a particular interpretation to be valid.  The second kind is constructive:
signed routing, affine deletion, persistent pair tracks, sparse execution, and optional MSA
supervision are possible ways to meet or approximate those requirements.  The constructions are
examples, not definitions of the framework.

Two limitations are especially revealing.  A deterministic graph router that must remain globally
invariant to all active endpoint identities can collapse to a fixed graph.  Relaxing the claim to a
targetwise directed conditional operator avoids that collapse but gives up the semantics of one
globally identified undirected graph.  Separately, a router forbidden from observing two endpoint
residues faces an irreducible Bayes error whenever the desired edge label depends on information
contained only in those residues.  Training data and teacher supervision can approach that floor;
they cannot remove it at inference.

Realizability is therefore not an afterthought.  It marks the boundary between a mathematically
well-defined interaction and a trainable mechanism that can infer or execute it under finite
information and compute.
